% 用户更正后已审核冻结 Appendix D。以下为 4.7 与 Appendix D 的定点替换候选；
% 已接受 2.7 未作为本任务冻结输入，符号接口仍须由父代理核对。

% 替换 sec-4-7 第77--94行：
可用弱馈送与后续退激相结合的条件概率模型描述缓发中子发射。令 $I_\beta(i)$ 为每次母核总衰变中直接馈送子核状态 $i$ 的概率，包括连续能谱时求和含积分，且 $\sum_i I_\beta(i)=b_\beta\leq1$。选择要计入的退激过程和观测时窗，对任意非终止状态 $a$ 定义
\begin{equation}\label{eq:4-7-cascade-pi-def}
\Pi_x(a)=\sum_{b,r}\frac{\Gamma^{(r)}_{a\to b}}{\Gamma_{\rm tot}(a)}\,\Pi_{x-r}(b),\qquad \Pi_{x<0}(a)=0.
\end{equation}
这里 $r$ 是这一步释放的中子数；$\gamma$ 跃迁有 $r=0$，单中子跃迁有 $r=1$。其他粒子发射或裂变通道若开放，亦应计入总宽度和相应后继状态或终止类别；若需区分剩余核与裂变，可将终止类别作为概率的附加指标。对于每一个所选吸收终态 $t$，规定 $\Pi_x(t)=\delta_{x0}$。当所有开放通道均被计入且体系以概率 1 在所定义时窗或理想终止条件下吸收时，$\sum_{x\ge0}\Pi_x(a)=1$。

按每次母核总衰变归一的 $\beta$ 缓发 $x$ 中子分支与按 $\beta$ 衰变条件归一的概率分别为
\begin{equation}\label{eq:4-7-conditional-probability}
P_{xn}^{(\beta;{\rm parent})}=\sum_i I_\beta(i)\Pi_x(i),\qquad
P_{xn}^{(\beta\mid\beta)}=\frac{P_{xn}^{(\beta;{\rm parent})}}{b_\beta}\quad(b_\beta>0).
\end{equation}
前者对全部 $x$ 求和为 $b_\beta$，后者为 1。至少一个缓发中子的概率及其平均多重度分别为 $\sum_{x\ge1}P_{xn}^{(\beta;{\rm parent})}$ 和 $\sum_{x\ge1}xP_{xn}^{(\beta;{\rm parent})}$；改变归一口径时应同时改变两式。仅当后续中子发射恰为一次且竞争 $\gamma$ 路径不会再产生中子时，单步近似 $\Pi_1(i)\simeq\Gamma_n(i)/\Gamma_{\rm tot}(i)$ 才成立。

% 第127行建议替换：
令 $R_{y|x}^{(\beta)}$ 为 $\beta$ 级联实际发射 $x$ 个中子时记录 $y$ 个信号的条件概率；各分布需按中子能谱、门及实验时窗加权。若按每次母核总衰变归一，则 $\beta$ 事件对观测的贡献为 $p_y^{(\beta)}=\sum_x R_{y|x}^{(\beta)}P_{xn}^{(\beta;{\rm parent})}$，其总和为 $b_\beta$。无条件母核数据还须加入所有非 $\beta$ 通道的贡献；其中无中子通道可贡献零中子事件，而自发裂变等也可能产生中子。若使用 $\beta$ 门，则应采用条件归一与门效率相应修正。

% 第30--44行保留公式，并在其后增加：
式\eqref{eq:4-7-differential-spectrum}为所述允许近似下的最低阶写法；辐射修正、有限核尺寸与反冲等需按所采用的精度和约定加入。文献若将 $g_A/g_V$ 吸收入 $B(GT)$ 或其积分常数 $K$，比较时须先统一定义，不再重复乘入。式中的 $F$ 与无量纲化的形状因子均为无量纲量。

% 第47--51行中“保护仅针对零动量转移下的总矢量荷”改为：
对于这里讨论的 Fermi 强度，直接使用的是同位旋对称极限下零动量总矢量荷的归一；CVC 也对矢量流及相关形状因子之间施加关系，并非对有限动量和截断模型中的全部矢量矩阵元逐一保证不变。

% 第57行的 Wang 结果改为范围化表述：
在 Wang 等所示的特定 SkM* 计算中，PVC 将 $^{56}\mathrm{Ni}$ 及 $^{132}\mathrm{Sn}$ 的部分 $1^+$ 强度移入相应衰变窗口；$^{56}\mathrm{Ni}$ 的所引低能跃迁涉及电子俘获。该结果随所选能量泛函及相关模型变化，例如作者所列 SGII 的 $^{132}\mathrm{Sn}$ PVC 计算未在窗口内得到 GT 强度。

% 第131--141行公式本身可保留，在式(4-7-NME-decomposition)后补充：
这里给出的分解是长程贡献；在所采用的轻中微子交换有效理论与阶次下，总算符还应包括相容归一的短程接触项，故 $M_{0\nu}$ 为相应总振幅。外提 $g_A^4$ 后，分解中的 $g_V/g_A$、电流形式因子及短程项仍可能随所采用的耦合和归一约定变化；不能仅凭外提因子推断改变 $g_A$ 时半衰期必严格按 $g_A^{-4}$ 缩放。第156行关于领头阶接触项的表述应保留第141行规定的轻中微子交换与非微扰两核子散射框架。

% 替换 Appendix D 第134--140行：
取单核子同位旋约定 $t_3|n\rangle=\frac12|n\rangle$、$t_3|p\rangle=-\frac12|p\rangle$，以及 $t_-|n\rangle=|p\rangle$、$t_+=t_-^\dagger$，故 $[t_+,t_-]=2t_3$。定义裸一体算符 $\mathbf O_-=\sum_k\boldsymbol\sigma(k)t_-(k)$、$\mathbf O_+=\mathbf O_-^\dagger$；对初态 $J_i$ 的磁子态取等权平均，并对两个电荷道各自全部末态及自旋分量求和：
\[
S_\mp=\frac{1}{2J_i+1}\sum_{M_i,f,M_f,\mu}
|\langle fJ_fM_f|O_{\mp,\mu}|iJ_iM_i\rangle|^2
=\sum_f\frac{|\langle fJ_f\|\mathbf O_\mp\|iJ_i\rangle|^2}{2J_i+1}.
\]
在两支末态完备、裸一体算符及上述约定下，由闭合与对易关系得到 $S_--S_+=3(N-Z)$。若只以常数 $q$ 缩放两个算符，差和变为 $q^2 3(N-Z)$；若加入两体流或截断空间有效算符，应重新求其对易子与闭合关系，不能不加条件沿用裸一体差和。有限 $Q_\beta$ 窗中的强度与完整差和是不同量。


% Appendix D 第199--234行：
为了保留经过多步退激后释放的中子总数，可把瞬态写为增广状态 $(a,m)$，其中 $a$ 标记当前核素和能级，$m$ 是自初始馈送起已经发射的中子数。对于发射 $r$ 个中子的单步过程，转移为 $(a,m)\to(b,m+r)$，概率为 $\Gamma_{a\to b}^{(r)}/\Gamma_{\rm tot}(a)$。对所选终态类别 $c$，吸收态写为 $(c,x)$。在有限状态、几乎必然吸收且 $\rho(Q)<1$ 的条件下，
\[
\mathbf N=(\mathbf I-\mathbf Q)^{-1},\qquad
\mathbf B=\mathbf N\mathbf R,\qquad
\Pi_{x,c}(a)=B_{(a,0),(c,x)},\quad
\Pi_x(a)=\sum_c\Pi_{x,c}(a).
\]
也可不显式扩维，而用式\eqref{eq:4-7-cascade-pi-def}对剩余中子数递推。连续谱情形需将求和推广为积分核；观测时窗内的长寿命态是否作为终态由实验定义。只有覆盖全部通道且最终吸收时，全多重度归一才成立。

% Appendix D 第56--70行：保留最终式( appD-seniority-e2-scaling )，替换其推导：
这里参考态 $n=v$ 是准自旋最低权重态，其投影为 $S_0=-S$。对于 $S>0$，正确的 CG 比是
\[
\frac{\langle S,S_0;1,0\mid S,S_0\rangle}
     {\langle S,-S;1,0\mid S,-S\rangle}
=\frac{S_0}{-S}=\frac{\Omega_j-n}{\Omega_j-v}.
\]
其中 $0\le v\le\Omega_j$，$n=v,v+2,\ldots,2\Omega_j-v$，且须存在相应资数态；只有当该阶梯实际包含 $n=\Omega_j$ 时才能取其半满点。$v=\Omega_j$ 时 $S=0$，分式不能使用；准自旋矢量在两个 $S=0$ 态之间的矩阵元由三角条件单独判为零。

% Appendix D 第174--188行：
在轻 Majorana 中微子交换机制下，两个弱流顶点由一条虚中微子传播子连接，传播子积分覆盖宽广动量范围。对所采用的 $g_A$、相空间及核半径归一，应将长程矩阵元与所需短程接触项一致地组合为总矩阵元。Cirigliano 等在轻中微子交换、非微扰 $^1S_0$ 两核子散射的领头阶重整化分析中引入短程项；坐标空间的接触作用在具体计算中需与所选正则化及耦合常数配套，例如
\[
V_{\nu,\mathrm{short}}^{(R)}(\mathbf r_{12})
=-2g_\nu^{NN}(R)\,\tau_{1,+}\tau_{2,+}\,\delta_R^{(3)}(\mathbf r_{12}).
\]
其有限耦合可通过相容的格点 QCD 匹配确定；电磁电荷独立性破缺数据在作者的分析中约束有关耦合的组合，并可在附加假设下提供估计，不能单凭该数据无条件唯一确定 $g_\nu^{NN}$。上述系数的相位和 $\tau_\pm$ 定义应与引用文献的同位旋约定一起转换。

% Appendix D 第257--261行：
若 $P_l$ 已定义为中子道的 R 矩阵穿透因子，则在简化的单道、非带电粒子与平滑约化宽度条件下 $\Gamma_n=2\gamma_l^2P_l$；阈值附近 $P_l\propto k^{2l+1}\propto(\Delta E)^{l+1/2}$。不能在这个定义下再乘一个 $\sqrt{\Delta E}$；若采用不同定义的势垒透射系数，则应重新给出该系数与相空间因子的关系。$\Gamma_n/\Gamma_{\rm tot}$ 的大小仍需实际约化宽度和竞争道输入，单凭阈值不产生固定的 $90\%$ 分支。

% Appendix D 第157--160行：
秩为 1 的项应把矢量流的 $\mathbf r$ 与 $\mathbf p/m_N$ 项和轴矢量流的 $[\boldsymbol\sigma\otimes\mathbf r]^{(1)}$ 项分别列出，不将 $\mathbf r$ 一并标成轴矢量算符。
